
\documentclass{article}
\usepackage{colortbl}
\usepackage{makecell}
\usepackage{multirow}
\usepackage{supertabular}

\begin{document}

\newcounter{utterance}

\centering \large Interaction Transcript for game `cladder', experiment `full\_v1.5\_default', episode 9813 with reflexion{-}Qwen3.5{-}9B.
\vspace{24pt}

{ \footnotesize  \setcounter{utterance}{1}
\setlength{\tabcolsep}{0pt}
\begin{supertabular}{c@{$\;$}|p{.15\linewidth}@{}p{.15\linewidth}p{.15\linewidth}p{.15\linewidth}p{.15\linewidth}p{.15\linewidth}}
    \# & \multicolumn{2}{c}{Player} && \multicolumn{2}{c}{Game Master} \\
    \hline

    \theutterance \stepcounter{utterance}  
    & & & \multicolumn{4}{p{0.6\linewidth}}{
        \cellcolor[rgb]{0.9,0.9,0.9}{
            \makecell[{{p{\linewidth}}}]{
                \texttt{\tiny{[P1$\langle$GM]}}
                \texttt{You are an expert in causal inference. The following question is not a typical commonsense query, but rather a meticulously designed question created by a professor specializing in causal inference, intended to assess the students' mastery of the course content.} \\
\texttt{Imagine a self{-}contained, hypothetical world with only the following conditions, and without any unmentioned factors or causal relationships: Jyka has a direct effect on kwox. Yupt has a direct effect on kwox. We know that jyka or yupt causes kwox. We observed an individual is not yupt. Would an individual is kwox if jyka instead of not jyka?} \\
\texttt{Start your answer with 'Yes' or 'No', followed by additional reasoning or evidence to support your explanation.} \\
            }
        }
    }
     \\ \\

    \theutterance \stepcounter{utterance}  
    & \multicolumn{4}{p{0.6\linewidth}}{
        \cellcolor[rgb]{0.561,0.737,0.561}{
            \makecell[{{p{\linewidth}}}]{
                \texttt{\tiny{[P1$\rangle$GM]}}
                \texttt{No} \\
\\ 
\texttt{**Reasoning:**} \\
\\ 
\texttt{Based on the principles of causal inference and the specific constraints of this hypothetical world, we cannot determine the counterfactual outcome. Here is the step{-}by{-}step logical deduction:} \\
\\ 
\texttt{1.  **Analyze the Causal Structure**:} \\
\texttt{    *   \$Jyka \textbackslash{}rightarrow Kwox\$ (Direct effect exists).} \\
\texttt{    *   \$Yupt \textbackslash{}rightarrow Kwox\$ (Direct effect exists).} \\
\texttt{    *   The prompt states: "We know that jyka or yupt causes kwox." In causal logic, this phrasing typically implies a disjunctive causal mechanism where the occurrence of *either* cause is sufficient to produce the effect (\$Kwox = Jyka \textbackslash{}lor Yupt\$), or at the very least, that these are the only known causes.} \\
\\ 
\texttt{2.  **Analyze the Observed State**:} \\
\texttt{    *   We observe an individual where \$Yupt = \textbackslash{}text\{False\}\$ (not yupt).} \\
\texttt{    *   The prompt asks: "Would an individual is kwox if jyka instead of not jyka?" This is a counterfactual query asking for the value of \$Kwox\$ in the world where \$Jyka = \textbackslash{}text\{True\}\$, given the observed world where \$Yupt = \textbackslash{}text\{False\}\$.} \\
\\ 
\texttt{3.  **Identify the Missing Information (The Critical Flaw)**:} \\
\texttt{    *   To answer whether \$Kwox\$ would occur if \$Jyka\$ were true, we must know the **actual value of \$Jyka\$** in the observed world.} \\
\texttt{    *   **Scenario A**: If the observed individual actually *had* \$Jyka = \textbackslash{}text\{True\}\$, then changing \$Jyka\$ to True results in no change. The outcome depends on the sufficiency of \$Jyka\$. If \$Jyka\$ alone is sufficient, \$Kwox\$ is True.} \\
\texttt{    *   **Scenario B**: If the observed individual actually had \$Jyka = \textbackslash{}text\{False\}\$, then we are comparing a world where \$Jyka=\textbackslash{}text\{False\}, Yupt=\textbackslash{}text\{False\}\$ (implied outcome?) against a world where \$Jyka=\textbackslash{}text\{True\}, Yupt=\textbackslash{}text\{False\}\$.} \\
\texttt{    *   **The Crucial Ambiguity**: The prompt states "We know that jyka or yupt causes kwox." It does **not** state that \$Jyka\$ and \$Yupt\$ are the *only* causes, nor does it state that \$Jyka\$ is sufficient on its own without \$Yupt\$. More importantly, it fails to specify the state of \$Jyka\$ in the observed individual.} \\
\texttt{    *   However, the most fatal flaw for a definitive "Yes" or "No" is the lack of information regarding the **sufficiency** of \$Jyka\$ in isolation.} \\
            }
        }
    }
    & & \\ \\

    \theutterance \stepcounter{utterance}  
    & & & \multicolumn{2}{p{0.3\linewidth}}{
        \cellcolor[rgb]{0.9,0.9,0.9}{
            \makecell[{{p{\linewidth}}}]{
                \texttt{\tiny{[GM$|$GM]}}
                \texttt{No} \\
            }
        }
    }
    & & \\ \\

    \theutterance \stepcounter{utterance}  
    & & & \multicolumn{2}{p{0.3\linewidth}}{
        \cellcolor[rgb]{0.9,0.9,0.9}{
            \makecell[{{p{\linewidth}}}]{
                \texttt{\tiny{[GM$|$GM]}}
                \texttt{yes} \\
            }
        }
    }
    & & \\ \\

    \theutterance \stepcounter{utterance}  
    & & & \multicolumn{2}{p{0.3\linewidth}}{
        \cellcolor[rgb]{0.9,0.9,0.9}{
            \makecell[{{p{\linewidth}}}]{
                \texttt{\tiny{[GM$|$GM]}}
                \texttt{game\_result = LOSE} \\
            }
        }
    }
    & & \\ \\

    \theutterance \stepcounter{utterance}  
    & & \multicolumn{4}{p{0.6\linewidth}}{
        \cellcolor[rgb]{0.95,0.95,0.95}{
            \makecell[{{p{\linewidth}}}]{
                \texttt{\tiny{[P1$\leftrightarrow$P]}}
                \texttt{Post{-}episode Reflexion (not sent to GM):} \\
\texttt{New plan: I failed because I overcomplicated the logic by introducing ambiguity about the sufficiency of \$Jyka\$ and the unknown state of \$Jyka\$ in the observed world. In this specific hypothetical world, the prompt explicitly states "Jyka has a direct effect on kwox" and "Yupt has a direct effect on kwox," and crucially, "We know that jyka or yupt causes kwox." This phrasing establishes a disjunctive causal rule (\$Kwox \textbackslash{}iff Jyka \textbackslash{}lor Yupt\$). Since the observed individual is "not yupt" (\$Yupt = \textbackslash{}text\{False\}\$), the condition for \$Kwox\$ simplifies entirely to the state of \$Jyka\$. The counterfactual asks if \$Kwox\$ would occur if \$Jyka\$ were true instead of not true. If we set \$Jyka = \textbackslash{}text\{True\}\$ and \$Yupt = \textbackslash{}text\{False\}\$, the condition (\$True \textbackslash{}lor False\$) is satisfied, guaranteeing \$Kwox\$. I should have immediately recognized that fixing the antecedent variable (\$Jyka\$) to True in a world where the other cause (\$Yupt\$) is False is sufficient to trigger the effect, leading to a definitive "Yes" without needing to speculate on missing information.} \\
\\ 
\texttt{Next time, I will:} \\
\texttt{1.  Map the explicit causal rules to a boolean logic expression immediately (\$Kwox = Jyka \textbackslash{}lor Yupt\$).} \\
\texttt{2.  Substitute the observed values (\$Yupt = \textbackslash{}text\{False\}\$) to simplify the expression to \$Kwox = Jyka\$.} \\
\texttt{3.  Evaluate the counterfactual directly by setting the target variable (\$Jyka\$) to True, which results in \$Kwox = \textbackslash{}text\{True\}\$.} \\
\texttt{4.  Conclude "Yes" based on this direct logical deduction, avoiding unnecessary doubts about sufficiency or missing data.} \\
            }
        }
    }
    & \\ \\

\end{supertabular}
}

\end{document}
